\subsubsection{移动网格独立验证}

另在物理坐标中建立移动控制体模型，以相同物性、环境、时间步长和首达判据独立计算，与固定材料坐标解比较，结果见\cref{tab:q4movingvalidation}和\cref{fig:q4movingvalidation}。

\begin{table}[H]
\centering
\caption{移动网格与固定域映射的数值比较}\label{tab:q4movingvalidation}
\begin{tabular}{lrrr}
\toprule
指标 & 固定域映射法 & 移动网格法 & 相对差\\
\midrule
$t_d$/h & 50.821334575 & 50.821334575 & $5.39\times10^{-13}$\\
6 h，中心 & 1.716433965 & 1.716433965 & $1.95\times10^{-14}$\\
6 h，动态表面 & 0.421517450 & 0.421517450 & $3.59\times10^{-13}$\\
24 h，中心 & 0.283703339 & 0.283703339 & $1.54\times10^{-13}$\\
24 h，动态表面 & 0.067090805 & 0.067090805 & $3.87\times10^{-14}$\\
48 h，中心 & 0.155610914 & 0.155610914 & $2.58\times10^{-13}$\\
48 h，动态表面 & 0.052857529 & 0.052857529 & $1.06\times10^{-13}$\\
$t_d^{\rm map}$，中心 & 0.150000000 & 0.150000000 & $3.40\times10^{-13}$\\
$t_d^{\rm map}$，动态表面 & 0.052496307 & 0.052496307 & $9.44\times10^{-14}$\\
\bottomrule
\end{tabular}
\par\smallskip
{\footnotesize 水分单位为 kg/kg；相对差以固定域结果为分母。末两行为同一固定域达标时刻的比较。}
\end{table}


两法达标时间相对差为 $5.81\times10^{-13}$，全水分场最大差为 $2.61\times10^{-13}$~kg/kg、温度场最大差为 $1.16\times10^{-10}~{}^\circ\mathrm C$，并通过几何守恒检查。

\begin{figure}[H]
\centering
\includegraphics[width=0.96\textwidth]{q4_moving_mesh_validation.pdf}
\caption{两种坐标实现的中心水分轨迹及物理半径剖面对比；曲线在图示精度下重合，圆点仅用于区分移动网格曲线}\label{fig:q4movingvalidation}
\end{figure}

两种坐标实现的差异处于浮点舍入量级，支持程序实现的一致性；该比较不替代对收缩假设的实验验证。
